% SHARD-ID: AQM-17-PROJECTIVE-LINE-STALKS
% SHARD-TITLE: Projective Line Stalks
% SHARD-SUMMARY: Defines stalks as germs and applies the Proj stalk formula to \(\mathbf P^1_{\mathbb F_3}\).
% SHARD-SUMMARY: Computes the local rings and \(\mathcal O(d)\) stalks at each rational point used in the projective-line example.
% SHARD-SUMMARY: Relates the stalk residues to the evaluation rows used by the finite symplectic field pairing.
% SHARD-KEYWORDS: stalk, germ, projective line, local ring, O(d), rational point

\section{Projective Line Stalks}
\label{sec:projective-line-stalks}

\subsection{Status and Local Evidence}

This shard continues Shard \(AQM\)-16.  The phrase ``each point in the
example'' means the four rational points used by the finite pairing:
\([1:0]\), \([0:1]\), \([1:1]\), and \([2:1]\).  The scheme has more points
than these four; a general formula for those other stalks is recorded at the
end.

The source definition of a stalk is local:
\path{references/algebraic_geometry/stacks_project_sheaves.tex}, lines
894--944.  The source formula for stalks on \(\operatorname{Proj}(S)\) is
\path{references/algebraic_geometry/stacks_project_constructions.tex}, lines
1091--1116 and 1182--1207.  The checked rows are produced by
\path{scripts/bridges/projective_line_sheaf_fields.jl}.

\subsection{What a Stalk Is}

Let \(X\) be a topological space, \(x\in X\), and \(\mathcal F\) a presheaf.
The stalk at \(x\) is
\[
  \mathcal F_x=\operatorname*{colim}_{x\in U}\mathcal F(U),
\]
where \(U\) runs over open neighbourhoods of \(x\).  Concretely, an element is
represented by \((U,s)\), with \(x\in U\) and \(s\in\mathcal F(U)\), and
\((U,s)\sim(U',s')\) when some open neighbourhood
\(U''\subset U\cap U'\) of \(x\) satisfies \(s|_{U''}=s'|_{U''}\).
Thus a stalk element is a germ: it remembers the section only arbitrarily near
the point.

\subsection{The Proj Stalk Formula}

Let
\[
  S=\mathbb F_3[X,Y],\qquad \deg X=\deg Y=1,\qquad
  X_{\rm sch}=\operatorname{Proj}(S)=\mathbf P^1_{\mathbb F_3}.
\]
A point of \(X_{\rm sch}\) is a homogeneous prime
\(\mathfrak p\subset S\) not containing the irrelevant ideal \((X,Y)\).
For a graded \(S\)-module \(M\), the associated sheaf \(\widetilde M\) has
stalk
\[
  (\widetilde M)_{\mathfrak p}=M_{(\mathfrak p)}.
\]
The right hand side consists of homogeneous fractions
\[
  m/f,\qquad m\in M,\ f\in S \text{ homogeneous},\quad
  \deg_M m=\deg_S f,\quad f\notin\mathfrak p.
\]
For a shifted module such as \(S(d)\), the equality of degrees is in the
shifted grading of \(S(d)\), while the denominator remains an element of
\(S\).
For the structure sheaf, this gives
\[
  \mathcal O_{X_{\rm sch},\mathfrak p}=S_{(\mathfrak p)}.
\]
For the twist \(\mathcal O(d)=\widetilde{S(d)}\),
\[
  \mathcal O(d)_{\mathfrak p}=S(d)_{(\mathfrak p)}.
\]

\subsection{The Two Standard Charts}

On the chart \(D_+(X)\), set \(u=Y/X\).  Then
\[
  D_+(X)=\operatorname{Spec}\mathbb F_3[u].
\]
On this chart we use the local frame \(e_X^{(d)}\) of \(\mathcal O(d)\)
corresponding to \(X^d\).  The section
\[
  X^{d-i}Y^i
\]
has local expression
\[
  u^i e_X^{(d)}.
\]

On the chart \(D_+(Y)\), set \(v=X/Y\).  Then
\[
  D_+(Y)=\operatorname{Spec}\mathbb F_3[v].
\]
On this chart we use the local frame \(e_Y^{(d)}\) corresponding to \(Y^d\).
The same section has local expression
\[
  v^{d-i}e_Y^{(d)}.
\]
On the overlap \(D_+(XY)\), \(u=1/v\), and the two frames satisfy
\[
  e_X^{(d)}=v^d e_Y^{(d)},\qquad e_Y^{(d)}=u^d e_X^{(d)}.
\]
This is why both local expressions describe the same germ where both charts
are available.

\subsection{Point P-Infinity: [1:0]}

The homogeneous prime is
\[
  \mathfrak p_\infty=(Y).
\]
Since \(X\notin\mathfrak p_\infty\), this point lies in \(D_+(X)\).  In the
coordinate \(u=Y/X\), the corresponding prime of \(\mathbb F_3[u]\) is
\[
  (u).
\]
Therefore
\[
  \mathcal O_{P_\infty}=\mathbb F_3[u]_{(u)},\qquad
  \mathfrak m_{P_\infty}=(u),\qquad
  \kappa(P_\infty)=\mathbb F_3[u]_{(u)}/(u)=\mathbb F_3.
\]
The \(\mathcal O(d)\)-stalk is free of rank one:
\[
  \mathcal O(d)_{P_\infty}
  =
  \mathbb F_3[u]_{(u)}e_X^{(d)}.
\]
For the monomial \(X^{d-i}Y^i\),
\[
  (X^{d-i}Y^i)_{P_\infty}=u^i e_X^{(d)}.
\]
Reducing modulo \((u)\) gives
\[
  u^i e_X^{(d)}
  \longmapsto
  \begin{cases}
  1\cdot \overline e_X^{(d)},& i=0,\\
  0,& i>0.
  \end{cases}
\]
This is exactly the evaluation rule used earlier at \([1:0]\).

\subsection{Point P0: [0:1]}

The homogeneous prime is
\[
  \mathfrak p_0=(X).
\]
Since \(Y\notin\mathfrak p_0\), use \(D_+(Y)\) and \(v=X/Y\).  The local
prime is \((v)\).  Hence
\[
  \mathcal O_{P_0}=\mathbb F_3[v]_{(v)},\qquad
  \mathfrak m_{P_0}=(v),\qquad
  \kappa(P_0)=\mathbb F_3.
\]
The twist stalk is
\[
  \mathcal O(d)_{P_0}=\mathbb F_3[v]_{(v)}e_Y^{(d)}.
\]
The monomial \(X^{d-i}Y^i\) has germ
\[
  v^{d-i}e_Y^{(d)}.
\]
Modulo \((v)\),
\[
  v^{d-i}e_Y^{(d)}
  \longmapsto
  \begin{cases}
  1\cdot \overline e_Y^{(d)},& i=d,\\
  0,& i<d.
  \end{cases}
\]

\subsection{Point P1: [1:1]}

The homogeneous prime is
\[
  \mathfrak p_1=(X-Y).
\]
On \(D_+(Y)\), it becomes
\[
  (v-1)\subset\mathbb F_3[v].
\]
Therefore
\[
  \mathcal O_{P_1}=\mathbb F_3[v]_{(v-1)},\qquad
  \mathfrak m_{P_1}=(v-1),\qquad
  \kappa(P_1)=\mathbb F_3.
\]
The twist stalk is
\[
  \mathcal O(d)_{P_1}=\mathbb F_3[v]_{(v-1)}e_Y^{(d)}.
\]
The germ of \(X^{d-i}Y^i\) is
\[
  v^{d-i}e_Y^{(d)}.
\]
Modulo \((v-1)\), we set \(v=1\), so
\[
  v^{d-i}e_Y^{(d)}
  \longmapsto
  1\cdot \overline e_Y^{(d)}
\]
for every \(i\).  This explains why every monomial evaluates to \(1\) at
\([1:1]\).

\subsection{Point P2: [2:1]}

The homogeneous prime is
\[
  \mathfrak p_2=(X-2Y)=(X+Y).
\]
On \(D_+(Y)\), it becomes
\[
  (v-2)\subset\mathbb F_3[v].
\]
Thus
\[
  \mathcal O_{P_2}=\mathbb F_3[v]_{(v-2)},\qquad
  \mathfrak m_{P_2}=(v-2),\qquad
  \kappa(P_2)=\mathbb F_3.
\]
The twist stalk is
\[
  \mathcal O(d)_{P_2}=\mathbb F_3[v]_{(v-2)}e_Y^{(d)}.
\]
The germ of \(X^{d-i}Y^i\) is
\[
  v^{d-i}e_Y^{(d)}.
\]
Modulo \((v-2)\), we set \(v=2=-1\).  Therefore
\[
  v^{d-i}e_Y^{(d)}
  \longmapsto
  2^{d-i}\overline e_Y^{(d)}.
\]
The residue is \(1\) when \(d-i\) is even and \(2\) when \(d-i\) is odd.

\subsection{Target-Valued Stalks}

The arithmetic field sheaf used in the example is
\[
  \mathcal O(d)\otimes_{\mathbb F_3}V,\qquad V=\mathbb F_3^2.
\]
At any rational point \(P\), its stalk is
\[
  (\mathcal O(d)\otimes V)_P
  =
  \mathcal O(d)_P\otimes_{\mathbb F_3}V.
\]
Because \(\mathcal O(d)_P\) is a free rank-one \(\mathcal O_P\)-module, this
is two copies of the local ring:
\[
  (\mathcal O(d)\otimes V)_P
  \cong
  \mathcal O_P q\oplus \mathcal O_P p.
\]
The residue at \(P\) is
\[
  (\mathcal O(d)\otimes V)_P
  \longrightarrow
  \kappa(P)\otimes_{\mathbb F_3}V
  =
  V,
\]
and this is the pointwise value used in the finite symplectic pairing.

\subsection{Other Scheme Points}

The finite field algebra in Shard \(AQM\)-16 uses only rational points, but
the scheme has more stalks.  The generic point has
\(\mathcal O_\eta=\mathbb F_3(v)\) and
\(\mathcal O(d)_\eta=\mathbb F_3(v)e_Y^{(d)}\).  A non-rational closed point
on \(D_+(Y)\) is given by an irreducible \(h(v)\in\mathbb F_3[v]\) of degree
\(>1\); then
\[
  \mathcal O_h=\mathbb F_3[v]_{(h)},\quad
  \mathcal O(d)_h=\mathbb F_3[v]_{(h)}e_Y^{(d)},\quad
  \kappa(h)=\mathbb F_3[v]/(h).
\]
These stalks are not sampled by the current rational-point pairing.
